\section{Additional Experimental Details}
\label{app:details}

\subsection{Benchmark Construction}
\label{app:benchmark}

Our benchmark draws from 12 established datasets, each contributing one domain. For each domain, we sample 60 normal and 60 anomalous images for the test split, totaling 1{,}440 test items. All sampling uses a fixed random seed (42) for reproducibility.

\paragraph{Domain Selection.}
Our initial data collection covered 14 domains. We excluded two: \textbf{Digital Forensics} because it requires detecting manipulation artifacts in compressed images, a signal-processing task rather than visual anomaly detection; and \textbf{Surveillance/Avenue} because it targets behavioral anomalies in video sequences, outside the scope of appearance-based single-image anomaly detection. The final benchmark comprises 12 domains encoded as D1--D12.

\paragraph{Source Datasets.}
\begin{itemize}[leftmargin=*,itemsep=2pt]
    \item \textbf{D1 (Industrial Manufacturing)}: MVTec-AD~\citep{bergmann2019mvtec}. 15 categories of manufactured products, sampled proportionally across all categories.
    \item \textbf{D2 (Retail Products)}: GoodsAD~\citep{liang2023goodsad}. Consumer products with physical damage, opening defects, and surface damage.
    \item \textbf{D3 (Complex Industrial)}: VisA~\citep{zou2022visa}. Selected for its multi-object categories (PCBs, capsules) and complex texture objects (candles, cashews, macaroni), testing robustness beyond single-object inspection.
    \item \textbf{D4 (Infrastructure)}: SDNET2018~\citep{dorafshan2018sdnet}. Concrete surface images (deck, pavement, wall) with crack annotations.
    \item \textbf{D5 (Logical Anomaly)}: MVTec-LOCO~\citep{bergmann2022mvtecloco}. Correctly assembled products vs.\ logical errors (wrong count, wrong position, wrong component type).
    \item \textbf{D6 (3D Product Inspection)}: Real3D-AD. Rendered RGB views of 3D point cloud scans of industrial products. Normal products vs.\ geometric deformations (bulges, sinks, contamination). Tests whether VLMs can detect shape anomalies from rendered 3D data.
    \item \textbf{D7 (Remote Sensing)}: LEVIR-CD+. Satellite imagery of buildings before and after change events. No-change (normal) vs.\ building construction, demolition, or damage. Tests cross-modality generalization to aerial imagery.
    \item \textbf{D8 (Dermatology)}: DermaMNIST~\citep{yang2023medmnist} and ISIC dermoscopic images. Melanocytic nevi (normal) vs.\ melanoma and other malignancies.
    \item \textbf{D9 (Brain MRI)}: BraTS2021~\citep{baid2021brats}. Axial FLAIR slices. Normal slices vs.\ slices containing glioma.
    \item \textbf{D10 (Liver CT)}: BMAD-Liver~\citep{bao2024bmad}. Axial CT slices of normal liver vs.\ focal lesions.
    \item \textbf{D11 (GI Endoscopy)}: HyperKvasir~\citep{borgli2020hyperkvasir}. Normal colon mucosa vs.\ polyps and other pathology.
    \item \textbf{D12 (Road Safety)}: BDD100K~\citep{yu2020bdd100k} (normal driving scenes) and RoadAnomaly21~\citep{lis2019detecting} (scenes with unexpected obstacles).
\end{itemize}

\paragraph{Reference Images.}
Each test item is paired with up to $k{=}4$ normal reference images retrieved using DINOv2 visual similarity from a pre-built per-domain index.

\subsection{Domain Knowledge Specifications}
\label{app:domain_knowledge}

Each domain is equipped with a structured knowledge specification provided to the VLM agents. The specification contains: (1) a domain description, (2) what normal samples look like, (3) specific anomaly criteria, and (4) common false positives. An example for D1:

\begin{quote}
\small
\textbf{Domain}: Industrial Manufacturing (MVTec-AD). \\
\textbf{Normal}: Defect-free manufactured products. \\
\textbf{Anomaly criteria}: Surface defects (scratches, stains, discoloration); structural defects (cracks, breaks, missing parts); contamination (foreign objects, residue). \\
\textbf{Common false positives}: Lighting variation, slight rotation, normal manufacturing tolerance.
\end{quote}

\subsection{Debate Prompt Templates}
\label{app:prompts}

\paragraph{Anomaly Proposer (cold start).}
The Proposer receives: (1) $k$ normal reference images, (2) the query image, (3) domain knowledge text, and (4) expert evidence reports. It outputs a JSON object containing a \texttt{normal\_profile} (structural description of what the object should look like) and a \texttt{claims} array. Each claim includes: \texttt{id}, \texttt{category}, \texttt{location} (bounding box or quadrant), \texttt{appearance}, \texttt{evidence} (linking to expert findings), \texttt{analysis}, \texttt{confidence} (0--1), and \texttt{severity}.

\paragraph{Normality Advocate.}
The Advocate receives the same images plus the Proposer's claims. For each claim, it outputs: \texttt{refute\_confidence} (0--1, how confident it is the observation is NOT anomalous), \texttt{counter\_evidence} (specific reasons), and \texttt{likely\_cause} (benign attribution such as lighting, texture variation, compression artifacts).

\paragraph{Refinement rounds.}
In subsequent rounds, both agents receive only the \texttt{TBD} claims from the previous round along with the opposing agent's arguments. The Proposer may revise confidence levels; the Advocate may strengthen or withdraw challenges.

\subsection{Patch Expert Implementation}
\label{app:patch_expert}

\paragraph{Feature Extraction.}
DINOv2 ViT-S/14 (\texttt{vit\_small\_patch14\_dinov2.lvd142m}) from \texttt{timm}. Input: $518 \times 518$ pixels, producing $37 \times 37 = 1{,}369$ patches of 384 dimensions. All features are L2-normalized.

\paragraph{Anomaly Score.}
For each query patch, cosine distance to nearest reference patch (Eq.~\ref{eq:patch_distance}). Anomaly score = mean of top 1\% distances ($\approx$14 patches).

\paragraph{Text Report Generation.}
The report includes: (1) raw mean patch distance with annotation ``higher = more different''; (2) global CLS-token similarity; (3) qualitative assessment from sigmoid calibration $\sigma(20 \cdot (d - 0.18))$ binned into five levels.

\subsection{Bootstrap Confidence Intervals}
\label{app:bootstrap}

All confidence intervals: 10{,}000 bootstrap resamples, fixed seed. 95\% CIs computed over 10 domain-level AUROCs.

\tbd{[Bootstrap CI table to be populated after experiments.]}

\subsection{Per-Category Breakdown}
\label{app:per_category}

\tbd{[Per-category AUROC for D1 (MVTec-AD, 15 categories) and D3 (VisA, multi-object and complex categories) to verify gains are not driven by a subset.]}

\subsection{Detailed Failure Analysis}
\label{app:failure}

\tbd{[Detailed error analysis for worst-performing domain, including patch distance distributions, debate transcript examples, and VLM error patterns.]}
